\section{问题背景}
光伏发电能够为小区提供电力，但发电充足的时段与居民集中用电的时段未必一致。与外部电网相连的微网，可以从外网购电补充供给，也可以利用储能设备把电量留到以后使用。当光伏有剩余或外网电价较低时充电，在供电不足或电价较高时放电，是协调本地发电与用电需求的一种方式。

这种调控同时受到供电可靠性与购电成本的约束。提前多买电可能造成浪费，买得太少又可能在实际运行中产生高价补购；电池虽然能缓冲供需变化，却受到储电容量、充放电功率和能量损耗的限制。因此，微网需要根据决策时已经掌握的信息，协调各时段的购电与储能使用，在满足小区负载的前提下尽可能减少费用。

\section{问题重述}
\subsection{研究对象与共同条件}
研究对象为由光伏、储能设备和外网共同供电的小区微网。附件1给出代表日的电价、负载及光伏预测，附件2给出2025年全年的实际负载和光伏功率，附件3给出每天分四次发布的光伏预报，附件4给出同期波动电价。附件5提供各问的结果文件格式。

储能设备最大容量为12000 kWh，运行储电量须保持在1200--10800 kWh之间，最大充电和放电功率均为5000 kW。题目给定2025年1月1日零点储电量为6000 kWh。功率以kW表示，购电量、充放电量及储电量以kWh表示；各问均须保证微网供电满足小区负载。

\subsection{四问的具体任务}
\textbf{问题一：给定一天供需与电价时的购电安排。}根据附件1和储能参数，在零点安排当天的外网购电及储能充放电，要求当天零点与24点储电量相等。题目中的“每天相同”是指每日重复相同的电价和负载曲线，曲线在一天内仍可随时段变化。需给出满足供电要求、尽可能节省费用的具体安排。

\textbf{问题二：实际供需变化下的每日购电计划。}保持附件1的日内电价，利用附件2所对应的实际运行情景，在每天零点制定当天计划。普通购电按计划量计费；实际运行中供给不足的部分通过紧急购电补足，价格为相应电价的5倍。需同时给出每日计划、实际储能使用及紧急购电情况。

\textbf{问题三：利用更新预报调整购电计划。}在零点取得未来24小时整点光伏预报后制定当天计划，并可在6点、12点和18点取得新预报后调整尚未执行的购电安排。增加的购电量按相应电价的1.5倍收费，减少的部分涉及50\%的违约电价。需给出计划及调整结果，并判断引入日内预报是否值得。

\textbf{问题四：波动电价下的两类购电安排。}使用附件4的电价，分别重新研究问题二和问题三。两条分支分别保留原有的信息条件与调整权限，比较波动电价下的购电、储能使用和费用表现。

\subsection{结果范围与呈现要求}
完整购电结果按10分钟间隔提交。问题一至问题三分别填写\texttt{result1.xlsx}、\texttt{result2.xlsx}和\texttt{result3.xlsx}；问题四的两条分支分别填写\texttt{result4-2.xlsx}和\texttt{result4-3.xlsx}。问题二至问题四的完整结果覆盖2025年2月1日至12月31日，除计划购电和储能记录外，分别保留紧急购电及题目要求的调整记录。

论文按题目指定格式展示重点时段和日期，完整文件保留全部要求的记录。购电展示包括10、12、14、16、18和20点起始的六个10分钟区间，以及全天购电量和费用；储能展示按六个4小时区间分别汇总充电和放电，并列出日初与日末储电量。问题二、三重点展示3月20日、6月21日、9月23日和12月21日，并列出这些日期的紧急购电情况；问题四给出两条分支的相应结果。

\section{问题分析}
\subsection{问题一：电量在不同时段之间如何分配}
本问的主要矛盾是供需和电价在时间上的错配。某个时段光伏较多，并不意味着另一个时段无需购电；提前充电虽可减少后续高价购电，却会占用储能空间并产生损耗。因而，不能只比较一天的总发电量与总负载，也不能简单地在最低价时段尽量买满，而要同时考虑电量何时可得、何时需要，以及电池能否及时充入和放出。

日初与日末储电量相等，使当天的费用比较具有一致的起点和终点。若允许日末电量低于日初，就可能通过消耗初始库存减少当天购电，而把补回库存的费用留到以后。首尾相等消除了这部分表面节省，但并不要求电池全天保持同一电量；日内充放电仍是调控的主要手段。

\subsection{问题二：计划先制定，偏差后出现}
本问新增的困难是决策与实际信息出现的先后关系。附件2虽包含全年实际值，但每天零点尚未发生的负载和光伏不能直接用于当天计划。第二问没有另行提供当天供需预报，故需依靠此前已经发生的数据形成判断，并在运行时根据实际情况使用电池和紧急购电。

电池能够在实际光伏偏少或负载偏高时提供补充，也能吸收部分反向偏差，但其缓冲能力有限。即使当天总电量够用，缺电也可能先于富余电量出现；达到最低储电量或放电功率上限时，电池仍无法补足当前缺口。紧急购电的高价格使计划不能只追求普通购电量少，而应权衡预先多买的费用与不足时的补购风险。

连续运行还使每天的安排彼此关联：当天结束时留下多少电，会影响次日的应对能力。数据从1月1日开始，完整结果却从2月1日开始，二者之间的运行过程决定了2月初的可用储电量，不能直接把6000 kWh当作2月1日的初态。


\subsection{问题三：新信息的收益能否覆盖修改代价}
新预报使尚未执行的购电安排有机会更贴近实际需求，但修改订单本身有成本。部分预测变化可以由储能吸收，此时不一定需要改购电计划；若预计缺口超出储能可承担的范围，提前调增才可能比事后紧急购电更合算。调减同样需要比较退回购电款与违约费用。因此，预报更新是重新判断的依据，并不意味着每次都应修改订单。

本问还需要区分“获得预报的价值”与“日内继续更新预报的价值”。直接比较问题二与问题三，会同时包含零点新增光伏预报和日内更新的影响。评价日内预报时，应在具有相同零点信息的条件下，比较保留零点安排与利用后续预报调整的表现，才能说明6点、12点和18点的新信息究竟带来什么作用。

\subsection{问题四：供需变化之外还要面对价格未知}
按照本文采用的信息条件，未来真实电价在决策时不可提前获知。电价波动由此既影响购电费用，也影响提前充电是否值得、应保留多少电量以及何时使用储能。沿用原购电轨迹再更换价格，只能得到该轨迹的重新计价，不能回答价格变化后应怎样重新安排。

本问的两条分支必须分别研究。问题二对应的分支仍在零点制定普通购电计划；问题三对应的分支可在既定预报时刻调整未来安排。只有保留这一区别，才能分清价格不确定性和日内调整能力各自带来的影响，避免把额外的调整权限误当作某种购电方法的收益。

\section{总体思路}
四问围绕同一供电系统展开，共享负载满足、光伏利用、储电量变化及设备容量和功率限制。研究首先从问题一的一天安排入手，理清购电与储能转移电量之间的关系；再在问题二中加入实际供需偏差、紧急补购和跨日衔接，使当天计划与实际执行相区分。问题一形成的购电方案仅用于问题二首日启动，后续日期的计划则依据逐日积累的信息制定。

问题三在这一基础上引入按发布时间获得的光伏预报及有代价的计划调整，重点考察提前纠正偏差的作用。问题四分别扩展问题二和问题三，将固定日内电价替换为决策时尚未知晓的波动电价。这样，四问可以共用供电与储能的基本关系，同时分别处理供需偏差、信息更新和价格不确定性；费用比较也可明确对应到普通购电、调整和紧急购电三个来源。
